\section{预备定义与记号}

\subsection{Fibonacci 数与 Zeckendorf 数字}

采用 Fibonacci 数列 $(F_k)_{k\ge 0}$：$F_0=0,F_1=1,F_{k+1}=F_k+F_{k-1}$。Zeckendorf 表示使用权重 $(F_k)_{k\ge 2}$，即序列
\[
1,2,3,5,8,13,21,34,55,\dots
\]
对任意整数 $n\ge 0$，存在唯一数字序列 $(z_k(n))_{k\ge 2}\in\{0,1\}^{\NN}$ 使得
\[
n=\sum_{k\ge 2} z_k(n)\,F_k,\qquad z_k(n)\,z_{k+1}(n)=0\ \ \forall k\ge 2.
\]
将 $(z_k(n))$ 称为 $n$ 的 Zeckendorf 数字。

\subsection{原始读出空间与稳定语言}

给定分辨率参数 $m\in\NN$，定义原始读出标签空间
\[
Y_m=\{0,1,\dots,2^m-1\}.
\]
定义长度为 $m$ 的稳定语言
\[
X_m=\{w\in\{0,1\}^m:\text{$w$ 不含子串 }11\}.
\]
写 $w=(w_m,\dots,w_1)$，其中 $w_m$ 为最高位，局部约束为 $w_{j+1}w_j=0$（$1\le j<m$）。

\subsection{距离与推前}

对有限集合上的概率分布 $\mu,\nu$，总变差距离定义为
\[
\norm{\mu-\nu}_{\TV}=\frac{1}{2}\sum_x\abs{\mu(x)-\nu(x)}.
\]
对确定性映射 $\Pi$ 的推前算子 $\Pi_*$ 定义为
\[
(\Pi_*\mu)(A)=\mu(\Pi^{-1}(A)),
\]
其中 $A$ 为目标空间的子集。

